% !TEX program = lualatex
\documentclass[11pt,a4paper]{article}
\usepackage{icat-common}

%-------------------------------------------------
%  独自コマンド・設定
%-------------------------------------------------

% セクション名をラテン語形式に変更

%-------------------------------------------------
\begin{document}

\title{知覚原理に基づくフラクタル則}
\author{千葉将臣}
\date{2026-04-19}
\maketitle

\begin{abstract}
本稿は、知覚原理から必然的に要請される構造的制約としての「フラクタル則」を定式化するものである。本則は、知覚が「ある」と「ない」の境界を生じる際の幾何学的必然であり、四句分別（Catuskoti）の各句と一対一の対応関係を持つ。
\end{abstract}

\section*{序言}
フラクタル構造は、知覚が無限の再帰分割を行う過程で生じる「認識の限界形状」である。知覚原理によって示された非対称な境界の生成は、以下の四つの法則を通じて、顕現（Aletheia）の幾何学的基盤を形成する。

\section{定義}

\begin{description}[style=multiline, leftmargin=3cm]
  \item[定義 2.1] \textbf{自己相似性}：任意の部分が全体と相似である性質。
  \item[定義 2.2] \textbf{尺度不変性}：観測の尺度を変更しても、構造的特徴が保存される性質。
  \item[定義 2.3] \textbf{収束}：無限の細分化プロセスが、特定の極限状態（$\dfix$）へと安定化する動的な性質。
  \item[定義 2.4] \textbf{非対称性}：知覚における「ある（是）」と「ない（非）」の次元的・質的な格差。
\end{description}

\section{公理}

\begin{description}[style=multiline, leftmargin=3cm]
  \item[公理 2.1] 知覚の再帰的二分は、空間的相似性を必然的に伴う。
  \item[公理 2.2] 二分の反復は、尺度に依存しない純粋な論理的操作である。
  \item[公理 2.3] 生成のプロセスは、知覚不可能な極限へと漸近する。
\end{description}

\section{定理：フラクタル四則}

\begin{description}[style=multiline, leftmargin=3cm]
  \item[定理 2.1（自己相似則 ↔ 是）]
  任意の部分が全体と相似であることは、「ある（是）」の自己同一的な顕現の必然である。
  
  \item[定理 2.2（尺度不変則 ↔ 非）]
  尺度を変えても同じ構造が現れることは、特定の絶対的尺度を持たない「ない（非）」の普遍的性質の幾何学的表現である。
  
  \item[定理 2.3（収束則 ↔ 亦是亦非）]
  無限の細分化が $\dfix$ へ収束することは、生成（発散）と安定（停止）が同時並行する矛盾状態の物理的帰結である。
  
  \item[定理 2.4（非対称則 ↔ 非是非非）]
  「ある」と「ない」の間に生じる知覚的非対称性は、系を立ち上げるメタレベルの特異点であり、顕現の原動力（生命力）である。
\end{description}

\section{補足}
これら四則は、知覚原理からフラクタル則へと至る変換経路を補完するものである。自己相似性と尺度不変性が「表現」の基盤を、収束則が「演算」の停止性を、非対称性が「発生」の非決定性をそれぞれ担保する。

\section{系}
したがって、フラクタル則の欠如は知覚の崩壊を意味する。我々が世界を整合的な構造として知覚し得るという事実は、直ちに系がこれら四つの幾何学的制約下で作動していることの証明となる。

\section{リーマン球面とメビウス変換によるトポロジカルな基礎付け}
前節までに示されたフラクタル四則は、認識の位相空間をリーマン球面 $\hat{\mathbb{C}} = \mathbb{C} \cup \{\infty\}$ 上の力学系としてモデル化することで、より厳密な幾何学的基礎付けを得る。

\subsection{メビウス変換のパラメータと四則の対応}
知覚による境界生成プロセス（推論の推移や時間発展）は、複素平面上の空間を操作するメビウス変換 $f(z) = \frac{az+b}{cz+d}$ （$ad-bc \neq 0$）の反復適用として数学的に記述される。この変換を構成する幾何学的性質は、フラクタル四則および四句分別と以下のように構造的な同型をなす。

\begin{itemize}
    \item \textbf{自己相似性（是）}：パラメータ $a$ および $d$（またはトレース $\tau = a+d$）が決定する空間の回転・拡大縮小作用。パターンの肯定的な維持を表す。
    \item \textbf{尺度不変性（非）}：行列式 $ad-bc$ に関し、パラメータ全体を定数倍しても変換の性質が不変であるという射影空間の性質。特定の絶対尺度の否定を表す。
    \item \textbf{収束則（亦是亦非）}：変換が持つ不動点（アトラクタおよびリペラ）に向かう軌道。引力と斥力が共存し、極限へと漸近する動的な摩擦を表す。
    \item \textbf{非対称性（非是非非）}：パラメータ $c \neq 0$ による空間の反転作用（$1/z$）。この反転操作が線形な対称性を破壊し、有限の局所空間を絶対的特異点である無限遠点 $\infty$ へと接続する。
\end{itemize}

\subsection{無限遠点 $\infty$ と極限集合の生成メカニズム}
本体系において、無限遠点 $\infty$ は、いかなる世俗的な尺度や計算も破綻する絶対的な観測の限界点（勝義諦・空）に対応する。複数の異なるメビウス変換からなる群 $G$ （クライン群等）を想定し、この無限遠点 $\infty$ を初期値として変換を無限回適用した軌跡の閉包を極限集合（Limit Set） $\Lambda$ と定義する。

$$ \Lambda = \overline{\{ g(\infty) \mid g \in G \}} $$

この定義式が示す通り、認識空間上に立ち現れるフラクタル（極限集合）とは、絶対的隔離空間である無限遠点 $\infty$ が、有限の論理構造（世俗のメビウス変換）によって等長円の内外へ反転・圧縮され続けた結果として生じる幾何学的な「計算の摩擦痕」に他ならない。
したがって、フラクタル則の四相は、到達不能な無限遠点を有限の認識系へと射影する際に不可避的に生じる、位相幾何学的な制約条件そのものであると結論付けられる。

\section{付録：四句分別マッピング}
\begin{itemize}
    \item \textbf{自己相似性} $\iff$ \textbf{是}（空間的確定・肯定）
    \item \textbf{尺度不変性} $\iff$ \textbf{非}（時間的未確定・否定）
    \item \textbf{収束則} $\iff$ \textbf{亦是亦非}（時空混淆・演算プロセス）
    \item \textbf{非対称性} $\iff$ \textbf{非是非非}（時空超越・メタ発生）
\end{itemize}

\end{document}

